% Native reconstruction of source slides 46--74.
\section[Inverse matrices]{The Inverse of a Square Matrix}

% Slide 46
\sectioncard{2.2}{The Inverse of a Square Matrix}

% Slide 47
\begin{frame}{From Division to Matrix Inverses}
\begin{itemize}
  \item Matrix multiplication is not commutative: $AB\ne BA$ in general.
  \item We do not define matrix division directly.
\end{itemize}
\medskip

For numbers, $6/2=3$ can be written $6\cdot2^{-1}=3$, and the reciprocal $2^{-1}$ satisfies
\[
2\cdot2^{-1}=1=2^{-1}\cdot2.
\]
This suggests replacing division by multiplication with an inverse.
\begin{alertblock}{Question}
Can a square matrix have a two-sided reciprocal?
\end{alertblock}
\end{frame}

% Slide 48
\begin{frame}{Invertible Matrices}
\begin{block}{Definition}
An $n\times n$ matrix $A$ is \textbf{nonsingular}, or \textbf{invertible}, if there is a matrix $B$ such that
\[
AB=I_n=BA.
\]
The matrix $B$ is the \textbf{inverse} of $A$.
\end{block}
\medskip
For example, check that
\[
\begin{bmatrix}1&2&3\\0&1&4\\0&0&1\end{bmatrix}^{-1}
=
\begin{bmatrix}1&-2&5\\0&1&-4\\0&0&1\end{bmatrix}.
\]
\end{frame}

% Slide 49
\begin{frame}{Existence and Uniqueness of the Inverse}
Two questions arise for every square matrix:
\begin{enumerate}
  \item Does an inverse exist?
  \item If it exists, is it unique?
\end{enumerate}
\begin{block}{Fact}
If $A$ is invertible, then its inverse is unique.
\end{block}
\textbf{Proof.} If $B$ and $C$ are both inverses of $A$, then
\[
B=BI=B(AC)=(BA)C=IC=C.
\]
We denote the unique inverse by $A^{-1}$, so
\[
AA^{-1}=A^{-1}A=I_n.
\]
\end{frame}

% Slide 50
\begin{frame}{Properties of Inverse Matrices}
\theorembox{Theorem 6}{Suppose $A$ and $B$ are invertible. Then}
\begin{enumerate}[a.]
  \item $A^{-1}$ is invertible and $(A^{-1})^{-1}=A$.
  \item $AB$ is invertible and $(AB)^{-1}=B^{-1}A^{-1}$.
  \item $A^T$ is invertible and $(A^T)^{-1}=(A^{-1})^T$.
\end{enumerate}
\medskip
\textbf{Proof of (b), one direction:}
\[
(AB)(B^{-1}A^{-1})
=A(BB^{-1})A^{-1}
=AIA^{-1}=I.
\]
Similarly, $(B^{-1}A^{-1})(AB)=I$.
\end{frame}

% Slide 51
\begin{frame}{Order Matters for Inverses}
For three or more invertible matrices,
\[
(ABC)^{-1}=C^{-1}B^{-1}A^{-1}.
\]
The order reverses.
\medskip

For the transpose property,
\[
A^T(A^{-1})^T=(A^{-1}A)^T=I,
\]
and
\[
(A^{-1})^TA^T=(AA^{-1})^T=I.
\]
Therefore
\[
(A^T)^{-1}=(A^{-1})^T.
\]
\end{frame}

% Slide 52
\begin{frame}{Cancellation with an Invertible Matrix}
For arbitrary matrices, the following implications can fail:
\[
\begin{array}{ll}
AB=0\centernot\implies A=0\text{ or }B=0,
&
A\ne0,\ B\ne0\centernot\implies AB\ne0,\\[1ex]
A^2=A\centernot\implies A=I\text{ or }A=0,
&
AB=AC,\ A\ne0\centernot\implies B=C.
\end{array}
\]
If $A$ is invertible, left cancellation becomes valid:
\[
\begin{array}{ll}
AB=0\implies B=0,
&
B\ne0\implies AB\ne0,\\[1ex]
A^2=A\implies A=I,
&
AB=AC\implies B=C.
\end{array}
\]
\end{frame}

% Slide 53
\begin{frame}[shrink=7]{Finding an Inverse from Unknown Entries}
Find the inverse of
\[
A=\begin{bmatrix}2&1\\-1&0\end{bmatrix}.
\]
Let $B=\begin{bmatrix}x_1&x_2\\x_3&x_4\end{bmatrix}$ and solve $AB=I$:
\[
\begin{bmatrix}2&1\\-1&0\end{bmatrix}
\begin{bmatrix}x_1&x_2\\x_3&x_4\end{bmatrix}
=\begin{bmatrix}1&0\\0&1\end{bmatrix}.
\]
Thus
\[
2x_1+x_3=1,\quad 2x_2+x_4=0,\quad -x_1=0,\quad -x_2=1,
\]
so
\[
B=\begin{bmatrix}0&-1\\1&2\end{bmatrix}.
\]
Both $AB=I$ and $BA=I$, hence $B=A^{-1}$.
\end{frame}

% Slide 54
\begin{frame}{A Nonzero Matrix May Be Singular}
Consider
\[
A=\begin{bmatrix}1&0\\0&0\end{bmatrix}.
\]
If $B=\begin{bmatrix}x_1&x_2\\x_3&x_4\end{bmatrix}$, then
\[
AB=\begin{bmatrix}x_1&x_2\\0&0\end{bmatrix}
\ne
\begin{bmatrix}1&0\\0&1\end{bmatrix}
\]
for every $B$. Thus $A$ is not invertible.
\begin{block}{Terminology}
An $n\times n$ matrix with no inverse is called \textbf{singular}（奇异的）.
\end{block}
Unlike numbers, nonzero matrices can be noninvertible.
\end{frame}

% Slide 55
\begin{frame}{Row Equivalence and Invertibility}
\theorembox{Recall: uniqueness of reduced echelon form}{Every matrix is row equivalent to one and only one reduced echelon matrix.}
\medskip

\theorembox{Theorem 7}{An $n\times n$ matrix $A$ is invertible if and only if $A$ is row equivalent to $I_n$. Any sequence of elementary row operations that reduces $A$ to $I_n$ transforms $I_n$ into $A^{-1}$.}
\end{frame}

% Slide 56
\begin{frame}{The Row-Reduction Algorithm for $A^{-1}$}
\begin{block}{Algorithm}
Place $A$ and $I_n$ side by side. Row reduce the augmented matrix:
\[
[\,A\mid I_n\,]\longrightarrow[\,I_n\mid A^{-1}\,].
\]
If the left half cannot be reduced to $I_n$, then $A$ is not invertible.
\end{block}
\medskip
\begin{exampleblock}{Example}
Find $A^{-1}$ for
\[
A=\begin{bmatrix}5&3\\1&1\end{bmatrix}.
\]
\end{exampleblock}
\end{frame}

% Slide 57
\begin{frame}{Row-Reduction Example}
\small
\[
\left[\begin{array}{cc|cc}5&3&1&0\\1&1&0&1\end{array}\right]
\sim
\left[\begin{array}{cc|cc}1&1&0&1\\0&-2&1&-5\end{array}\right]
\sim
\left[\begin{array}{cc|cc}1&1&0&1\\0&1&-\frac12&\frac52\end{array}\right]
\]
\[
\sim
\left[\begin{array}{cc|cc}
1&0&\frac12&-\frac32\\
0&1&-\frac12&\frac52
\end{array}\right].
\]
Therefore
\[
A^{-1}=
\begin{bmatrix}
\frac12&-\frac32\\[.3ex]
-\frac12&\frac52
\end{bmatrix}.
\]
\end{frame}

% Slide 58
\begin{frame}[shrink=7]{A $3\times3$ Inverse}
Find the inverse of
\[
A=\begin{bmatrix}1&2&3\\4&5&8\\3&4&6\end{bmatrix}.
\]
Row reduction gives
\[
\left[\begin{array}{ccc|ccc}
1&2&3&1&0&0\\
4&5&8&0&1&0\\
3&4&6&0&0&1
\end{array}\right]
\sim
\left[\begin{array}{ccc|ccc}
1&0&0&-2&0&1\\
0&1&0&0&-3&4\\
0&0&1&1&2&-3
\end{array}\right].
\]
Hence
\[
A^{-1}=\begin{bmatrix}-2&0&1\\0&-3&4\\1&2&-3\end{bmatrix}.
\]
\end{frame}

% Slide 59
\begin{frame}{Detecting a Singular Matrix}
Consider
\[
A=\begin{bmatrix}1&0&0\\0&1&2\\0&1&2\end{bmatrix}.
\]
Row reduction produces a zero row:
\[
\begin{bmatrix}1&0&0\\0&1&2\\0&1&2\end{bmatrix}
\sim
\begin{bmatrix}1&0&0\\0&1&2\\0&0&0\end{bmatrix}.
\]
Therefore $A$ cannot be row equivalent to $I_3$ and is not invertible.
\begin{block}{Remark}
A square matrix is singular if and only if its echelon form has a row of zeros.
\end{block}
\end{frame}

% Slide 60
\begin{frame}{Rank and Invertibility}
For an $n\times n$ matrix $A$,
\[
\boxed{A\text{ is invertible}\iff \rank(A)=n.}
\]
Equivalently,
\[
\boxed{A\text{ is singular}\iff \rank(A)<n.}
\]
An invertible square matrix has a pivot in every row and every column.
\end{frame}

% Slide 61
\subsection{Matrix Equations}
\begin{frame}{Matrix Equations}
The linear system
\[
\begin{cases}
a_{11}x_1+a_{12}x_2+\cdots+a_{1n}x_n=b_1,\\
a_{21}x_1+a_{22}x_2+\cdots+a_{2n}x_n=b_2,\\
\hfill\vdots\\
a_{n1}x_1+a_{n2}x_2+\cdots+a_{nn}x_n=b_n
\end{cases}
\]
has the matrix form
\[
A\vect x=\vect b,\qquad
A=(a_{ij}),\quad
\vect x=\begin{bmatrix}x_1\\\vdots\\x_n\end{bmatrix},\quad
\vect b=\begin{bmatrix}b_1\\\vdots\\b_n\end{bmatrix}.
\]
\end{frame}

% Slide 62
\begin{frame}{Solving $A\vect x=\vect b$ with an Inverse}
Suppose $A$ is invertible. Multiplying on the left by $A^{-1}$ gives
\[
A^{-1}A\vect x=A^{-1}\vect b
\quad\Longrightarrow\quad
\vect x=A^{-1}\vect b.
\]
If $\vect w$ is another solution, then
\[
A\vect w=\vect b
\quad\Longrightarrow\quad
\vect w=A^{-1}\vect b.
\]
Thus any solution must equal $A^{-1}\vect b$, proving uniqueness.
\end{frame}

% Slide 63
\begin{frame}{Unique Solutions from Invertibility}
\theorembox{Theorem 5}{If $A$ is an invertible $n\times n$ matrix, then for every $\vect b\in\R^n$, the equation $A\vect x=\vect b$ has the unique solution
\[
\vect x=A^{-1}\vect b.
\]}
\medskip
\begin{block}{Special case}
If $A$ is invertible, the homogeneous equation
\[
A\vect x=\vect0
\]
has only the trivial solution $\vect x=\vect0$.
\end{block}
\end{frame}

% Slide 64
\begin{frame}[shrink=7]{Solution Types of Linear Systems}
For a general $m\times n$ matrix $A$, the system $A\vect x=\vect b$ may have
\[
\text{no solution},\qquad \text{one solution},\qquad\text{or infinitely many solutions}.
\]
If $A$ is square and invertible, every $A\vect x=\vect b$ has exactly one solution:
\[
\vect x=A^{-1}\vect b.
\]
\medskip

The homogeneous system $A\vect x=\vect0$ always has at least the trivial solution. It has either
\[
\text{only }\vect x=\vect0
\qquad\text{or}\qquad
\text{infinitely many solutions}.
\]
If $A$ is invertible, only the trivial solution occurs.
\end{frame}

% Slide 65
\begin{frame}{Solving a $2\times2$ System with $A^{-1}$}
Solve
\[
-7x_1+3x_2=2,\qquad 5x_1-2x_2=1.
\]
Here
\[
A=\begin{bmatrix}-7&3\\5&-2\end{bmatrix},
\qquad
A^{-1}=\begin{bmatrix}2&3\\5&7\end{bmatrix},
\qquad
\vect b=\begin{bmatrix}2\\1\end{bmatrix}.
\]
Therefore
\[
\vect x=A^{-1}\vect b
=\begin{bmatrix}2&3\\5&7\end{bmatrix}
\begin{bmatrix}2\\1\end{bmatrix}
=\begin{bmatrix}7\\17\end{bmatrix}.
\]
\end{frame}

% Slide 66
\begin{frame}{Example: The Inverse-Matrix Algorithm}
Solve
\[
\begin{cases}
x_1+2x_2-x_3=-5,\\
2x_1-x_2+x_3=6,\\
x_1-x_2-3x_3=-3.
\end{cases}
\]
We will write the system as $A\vect x=\vect b$ and compute
\[
\vect x=A^{-1}\vect b.
\]
\end{frame}

% Slide 67
\begin{frame}{Example: Matrix Form}
For the preceding system,
\[
A=\begin{bmatrix}
1&2&-1\\
2&-1&1\\
1&-1&-3
\end{bmatrix},\qquad
\vect x=\begin{bmatrix}x_1\\x_2\\x_3\end{bmatrix},\qquad
\vect b=\begin{bmatrix}-5\\6\\-3\end{bmatrix}.
\]
The inverse is
\[
A^{-1}=\frac1{19}
\begin{bmatrix}
4&7&1\\
7&-2&-3\\
-1&3&-5
\end{bmatrix}.
\]
\end{frame}

% Slide 68
\begin{frame}{Example: Solution}
\[
\vect x=A^{-1}\vect b
=\frac1{19}
\begin{bmatrix}
4&7&1\\
7&-2&-3\\
-1&3&-5
\end{bmatrix}
\begin{bmatrix}-5\\6\\-3\end{bmatrix}
\]
\[
=\frac1{19}\begin{bmatrix}19\\-38\\38\end{bmatrix}
=\begin{bmatrix}1\\-2\\2\end{bmatrix}.
\]
Thus
\[
x_1=1,\qquad x_2=-2,\qquad x_3=2.
\]
\end{frame}

% Slide 69
\begin{frame}{General Matrix Equations $AX=B$}
Suppose $A$ is $n\times n$ and $X,B$ are $n\times s$. If $A$ is invertible,
\[
A^{-1}AX=A^{-1}B
\quad\Longrightarrow\quad
\boxed{X=A^{-1}B}.
\]
\begin{block}{Conditions for the inverse-matrix algorithm}
\begin{enumerate}
  \item The coefficient matrix $A$ is square.
  \item The coefficient matrix $A$ is invertible.
\end{enumerate}
\end{block}
\end{frame}

% Slide 70
\begin{frame}[shrink=10]{Three Common Matrix Equations}
\begin{columns}[T]
\column{.32\textwidth}
\centering
\[
AX=B
\]
\[
X=A^{-1}B
\]
\column{.32\textwidth}
\centering
\[
XA=B
\]
\[
X=BA^{-1}
\]
\column{.34\textwidth}
\centering
\[
AXB=C
\]
\[
X=A^{-1}CB^{-1}
\]
\end{columns}
\medskip
For example, solve $AX=B$ with
\[
A=\begin{bmatrix}0&-1&1\\1&0&0\\1&0&2\end{bmatrix},
\qquad
B=\begin{bmatrix}0&-1\\2&-5\\0&1\end{bmatrix}.
\]
Since
\[
A^{-1}=\begin{bmatrix}0&1&0\\-1&-\frac12&\frac12\\0&-\frac12&\frac12\end{bmatrix},
\]
\[
X=A^{-1}B=\begin{bmatrix}2&-5\\-1&4\\-1&3\end{bmatrix}.
\]
\end{frame}

% Slide 71
\begin{frame}[shrink=9]{Exercises: Matrix Equations I}
\textbf{1.} Solve $XA=B$ for
\[
A=\begin{bmatrix}1&0&-2\\2&4&-1\\0&1&1\end{bmatrix},
\qquad
B=\begin{bmatrix}1&0&2\\2&4&3\end{bmatrix}.
\]
Since
\[
A^{-1}=\begin{bmatrix}5&-2&8\\-2&1&-3\\2&-1&4\end{bmatrix},
\quad
X=BA^{-1}=\begin{bmatrix}9&-4&16\\8&-3&16\end{bmatrix}.
\]
\textbf{2.} Solve
\[
\begin{bmatrix}1&4\\0&1\end{bmatrix}
X
\begin{bmatrix}1&0\\3&1\end{bmatrix}
=\begin{bmatrix}-1&1\\1&0\end{bmatrix}.
\]
Multiplying by the two inverses gives
\[
X=\begin{bmatrix}-8&1\\1&0\end{bmatrix}.
\]
\end{frame}

% Slide 72
\begin{frame}[shrink=10]{Exercises: Matrix Equations II}
\textbf{3.} Solve $AX=A+2X$ for
\[
A=\begin{bmatrix}4&0&0\\0&1&0\\0&0&1\end{bmatrix}.
\]
Rearranging gives
\[
(A-2I)X=A,\qquad X=(A-2I)^{-1}A.
\]
\medskip

\textbf{4.} Solve $A^2-XA-I=0$, where
\[
A=\begin{bmatrix}1&0&3\\2&1&1\\-1&0&1\end{bmatrix},
\qquad
A^{-1}=\begin{bmatrix}.25&0&-.75\\-.75&1&1.25\\.25&0&.25\end{bmatrix}.
\]
Then
\[
X=(A^2-I)A^{-1}=A-A^{-1}
=\begin{bmatrix}.75&0&3.75\\2.75&0&-.25\\-1.25&0&.75\end{bmatrix}.
\]
\end{frame}

% Slide 73
\begin{frame}{Exercise: Isolating an Unknown Matrix}
Suppose $A,B,C$ are invertible $n\times n$ matrices and
\[
A=B(D-I_n)C.
\]
Multiplying on the left by $B^{-1}$ and on the right by $C^{-1}$ gives
\[
B^{-1}AC^{-1}=D-I_n.
\]
Therefore
\[
\boxed{D=B^{-1}AC^{-1}+I_n}.
\]
The order of the factors must be preserved.
\end{frame}

% Slide 74
\begin{frame}[shrink=9]{Exercise: A Matrix Equation with a Transpose}
Let
\[
ABA^T=2BA^T+I,\qquad
A=\begin{bmatrix}1&0&0\\0&1&2\\0&0&1\end{bmatrix}.
\]
Since $2BA^T=2IBA^T$,
\[
(A-2I)BA^T=I.
\]
Hence
\[
B=(A-2I)^{-1}(A^T)^{-1}
=\bigl[A^T(A-2I)\bigr]^{-1}
=(A^TA-2A^T)^{-1}.
\]
Now
\[
A^TA-2A^T=
\begin{bmatrix}-1&0&0\\0&-1&2\\0&-2&3\end{bmatrix},
\]
so
\[
\boxed{B=\begin{bmatrix}-1&0&0\\0&3&-2\\0&2&-1\end{bmatrix}}.
\]
\end{frame}
